\documentclass[11pt]{article}
\usepackage[utf8]{inputenc}
\usepackage[T1]{fontenc}
\usepackage{amsmath,amssymb}
\usepackage{graphicx}
\usepackage{booktabs}
\usepackage{hyperref}
\usepackage[margin=1in]{geometry}
\usepackage{natbib}

\title{A SUMO E3 Ligase Promotes Long Non-Coding RNA Transcription for Small RNA-Directed DNA Elimination in \textit{Tetrahymena}}

\author{
  Author~One\textsuperscript{1},
  Author~Two\textsuperscript{1},
  Author~Three\textsuperscript{1,2} \\
  \textsuperscript{1}Department of Molecular Biology, University \\
  \textsuperscript{2}Department of Cell Biology, University
}

\date{}

\begin{document}
\maketitle

% ============================================================
\begin{abstract}
Programmed DNA elimination in \textit{Tetrahymena thermophila} is guided by Piwi-associated small RNAs (scnRNAs) that distinguish germline-limited sequences from somatic sequences through pairing with nascent long non-coding RNAs (lncRNAs) transcribed from the somatic genome. The mechanism that actively drives this lncRNA transcription has remained unclear. Here we show that Ema2, a SUMO E3 ligase, is essential for somatic lncRNA accumulation and consequently for target-directed small RNA degradation (TDSD), DNA elimination, and production of viable sexual progeny. Ema2 interacts with the SUMO E2 conjugating enzyme Ubc9 in a RING domain-dependent manner and promotes SUMOylation of the transcription elongation factor Spt6. In Ema2 knockout cells, both Spt6 and RNA Polymerase~II show 3--4-fold reduced chromatin association at lncRNA-producing loci, and lncRNA levels are reduced 6--10-fold at peak conjugation. We propose a model in which Ema2-mediated SUMOylation of Spt6 facilitates the chromatin loading of the transcription machinery, thereby driving the lncRNA transcription that is a prerequisite for scnRNA-directed heterochromatin formation and DNA elimination.
\end{abstract}

% ============================================================
\section{Introduction}

Programmed genome rearrangement is a widespread phenomenon in ciliated protozoa, where it serves to shape the somatic macronucleus from the germline micronucleus during sexual development \citep{chalker2011,mochizuki2012}. In \textit{Tetrahymena thermophila}, approximately 6,000 internal eliminated sequences (IESs) are excised from the developing macronuclear genome, a process guided by scan RNAs (scnRNAs), a class of Piwi-associated small RNAs produced from the germline genome during meiosis \citep{mochizuki2002}.

The current model for DNA elimination involves a comparison mechanism: scnRNAs are initially produced from the entire germline genome, but those complementary to somatic sequences are selectively degraded through target-directed small RNA degradation (TDSD) upon pairing with nascent lncRNAs transcribed from the parental somatic macronucleus \citep{aronica2008,noto2015}. The surviving scnRNAs, enriched for germline-limited sequences, then guide heterochromatin formation and DNA elimination in the developing macronucleus.

While the downstream steps of scnRNA-guided heterochromatin formation have been characterized in considerable detail, the mechanism that actively promotes lncRNA transcription from the somatic genome has remained poorly understood. Given that this lncRNA transcription is essential for distinguishing somatic from germline sequences, its regulation is a critical gap in our understanding of DNA elimination.

SUMOylation, the covalent attachment of Small Ubiquitin-like Modifier (SUMO) proteins to target substrates, is a post-translational modification with well-established roles in transcription regulation across eukaryotes \citep{gareau2010,hendriks2017}. SUMOylation is catalyzed by an enzymatic cascade comprising an E1 activating enzyme, an E2 conjugating enzyme (Ubc9), and E3 ligases that provide substrate specificity. Whether SUMOylation plays a role in the lncRNA transcription that drives genome rearrangement in ciliates has not been explored.

Here we identify Ema2 as a SUMO E3 ligase that is essential for lncRNA accumulation during conjugation in \textit{Tetrahymena}. We demonstrate that Ema2 interacts with Ubc9, promotes SUMOylation of the transcription elongation factor Spt6, and facilitates the chromatin association of both Spt6 and RNA Polymerase~II at lncRNA-producing loci.

% ============================================================
\section{Related Work}

The scan RNA pathway in \textit{Tetrahymena} was first described by \citet{mochizuki2002}, who identified scnRNAs as meiosis-specific small RNAs associated with the Piwi protein Twi1. Subsequent work established the TDSD model, in which scnRNAs complementary to somatic sequences are selectively degraded \citep{aronica2008,noto2015}.

The transcription elongation factor Spt6 is a highly conserved histone chaperone that associates with RNA Polymerase~II and facilitates transcription through chromatin \citep{bortvin1996,ivanovska2011}. Spt6 has been shown to be subject to post-translational modifications that regulate its activity, but SUMOylation of Spt6 has not been previously described.

SUMOylation is increasingly recognized as a regulator of chromatin-associated processes, including transcription, DNA repair, and heterochromatin maintenance \citep{gareau2010,hendriks2017}. SUMO E3 ligases of the SP-RING family provide substrate specificity by bridging the E2 enzyme Ubc9 to target proteins \citep{johnson2001}. The connection between SUMOylation and programmed genome rearrangement has not been previously established.

% ============================================================
\section{Methods}

\subsection{Strains and culture conditions}

\textit{Tetrahymena thermophila} wild-type and Ema2 knockout strains were maintained in standard SPP medium. Conjugation was initiated by mixing starved cells of different mating types and monitored microscopically.

\subsection{Gene knockout and rescue}

Ema2 was disrupted by homologous recombination with a drug resistance cassette. Rescue strains were generated by reintroducing FLAG-tagged Ema2 under its endogenous promoter. RING domain mutants were constructed by site-directed mutagenesis.

\subsection{Progeny viability assays}

Sexual progeny viability was scored by isolating mating pairs at 24~hours post-mixing and testing for drug resistance indicative of successful macronuclear development. At least 100--150 pairs were scored per genotype.

\subsection{RT-qPCR for lncRNA quantification}

Total RNA was extracted at 0, 3, 6, 9, and 12~hours post-mixing. Somatic lncRNAs were quantified by RT-qPCR with locus-specific primers, normalized to rRNA levels.

\subsection{Co-immunoprecipitation}

Extracts from cells expressing FLAG-Ema2 and HA-Ubc9 were subjected to anti-FLAG immunoprecipitation followed by anti-HA western blotting to detect the Ema2--Ubc9 interaction. RING domain mutants were tested in parallel.

\subsection{In vivo SUMOylation assay}

Spt6 was immunoprecipitated from wild-type, Ema2-KO, and rescue strains, and SUMOylation was detected by anti-SUMO western blotting.

\subsection{Chromatin immunoprecipitation}

ChIP was performed at 8~hours post-mixing using antibodies against Spt6 and the Rpb1 subunit of RNA Polymerase~II. Enrichment at somatic lncRNA-producing loci was quantified by qPCR relative to IgG controls.

\subsection{Small RNA analysis and TDSD assay}

scnRNA levels were measured by northern blotting at 8~hours post-mixing. TDSD efficiency was assessed as the percentage of input scnRNAs remaining after the TDSD window.

% ============================================================
\section{Results}

\subsection{Ema2 is essential for DNA elimination and progeny viability}

To determine the function of Ema2 in sexual development, we crossed wild-type and Ema2 knockout strains and scored progeny viability (Table~\ref{tab:viability}). Wild-type crosses produced viable progeny in approximately 95\% of pairs, while Ema2 knockout crosses yielded only approximately 5\% viability. Reintroduction of Ema2 rescued viability to approximately 90\%, confirming that the phenotype is attributable to Ema2 loss.

\begin{table}[ht]
\centering
\caption{Sexual progeny viability in Ema2 mutants.}
\label{tab:viability}
\begin{tabular}{lccc}
\toprule
\textbf{Genotype} & \textbf{Pairs scored} & \textbf{Viable progeny (\%)} & \textbf{DNA elimination} \\
\midrule
WT $\times$ WT & 150 & $\sim$95 & Complete \\
Ema2-KO $\times$ Ema2-KO & 150 & $\sim$5 & Incomplete \\
Ema2-rescue $\times$ Ema2-rescue & 100 & $\sim$90 & Complete \\
\bottomrule
\end{tabular}
\end{table}

\subsection{Ema2 is required for somatic lncRNA accumulation}

We measured lncRNA levels from the somatic genome by RT-qPCR during conjugation (Table~\ref{tab:lncrna}). In wild-type cells, somatic lncRNAs accumulated dramatically, peaking at approximately 18.7-fold induction at 9~hours post-mixing. In Ema2 knockout cells, lncRNA accumulation was severely impaired, reaching only approximately 2.1-fold at the same time point---a reduction of approximately 9-fold relative to wild type.

\begin{table}[ht]
\centering
\caption{Somatic lncRNA accumulation during conjugation (fold change over 0~h).}
\label{tab:lncrna}
\begin{tabular}{lcc}
\toprule
\textbf{Time (h)} & \textbf{WT} & \textbf{Ema2-KO} \\
\midrule
0 & 1.0 & 1.0 \\
3 & 4.5 & 1.2 \\
6 & 12.3 & 1.8 \\
9 & 18.7 & 2.1 \\
12 & 15.2 & 1.9 \\
\bottomrule
\end{tabular}
\end{table}

\subsection{TDSD is impaired in Ema2 knockout cells}

Consistent with reduced lncRNA production, TDSD was severely impaired in Ema2 knockout cells. At 8~hours post-mixing, approximately 85\% of scnRNAs remained in Ema2-KO cells, compared with approximately 25\% in wild type (Table~\ref{tab:tdsd}). The failure of TDSD explains the downstream DNA elimination defect.

\begin{table}[ht]
\centering
\caption{TDSD efficiency at 8~hours post-mixing.}
\label{tab:tdsd}
\begin{tabular}{lcc}
\toprule
\textbf{Condition} & \textbf{scnRNA remaining (\%)} & \textbf{TDSD efficiency} \\
\midrule
WT & $\sim$25 & Normal \\
Ema2-KO & $\sim$85 & Severely impaired \\
\bottomrule
\end{tabular}
\end{table}

\subsection{Ema2 interacts with Ubc9 via its RING domain}

Co-immunoprecipitation experiments demonstrated that FLAG-Ema2 robustly co-precipitated HA-Ubc9 (Table~\ref{tab:coip}). This interaction was abolished by mutation of the RING domain, consistent with the RING-dependent E2 recruitment characteristic of SUMO E3 ligases. FLAG alone did not pull down HA-Ubc9.

\begin{table}[ht]
\centering
\caption{Co-immunoprecipitation of Ema2 and Ubc9.}
\label{tab:coip}
\begin{tabular}{lcc}
\toprule
\textbf{Bait} & \textbf{Prey detected} & \textbf{Signal} \\
\midrule
FLAG-Ema2 & HA-Ubc9 & +++ \\
FLAG-Ema2-RING-mut & HA-Ubc9 & -- \\
FLAG alone & HA-Ubc9 & -- \\
\bottomrule
\end{tabular}
\end{table}

\subsection{Ema2 promotes SUMOylation of Spt6}

In vivo SUMOylation assays revealed that Spt6 SUMOylation was reduced to approximately 15\% of wild-type levels in Ema2-KO cells (Table~\ref{tab:sumo}). Rescue with wild-type Ema2 restored SUMOylation to approximately 88\% of wild-type levels, establishing Ema2 as the major E3 ligase responsible for Spt6 SUMOylation.

\begin{table}[ht]
\centering
\caption{Spt6 SUMOylation levels.}
\label{tab:sumo}
\begin{tabular}{lc}
\toprule
\textbf{Condition} & \textbf{Spt6-SUMO intensity (AU)} \\
\midrule
WT & 1.00 \\
Ema2-KO & 0.15 \\
Ema2-rescue & 0.88 \\
\bottomrule
\end{tabular}
\end{table}

\subsection{Ema2 facilitates chromatin association of Spt6 and RNA Pol~II}

ChIP-qPCR at somatic lncRNA-producing loci demonstrated that both Spt6 and RNA Pol~II (Rpb1) chromatin occupancy were markedly reduced in Ema2-KO cells (Table~\ref{tab:chip}). Spt6 enrichment decreased from 8.5-fold to 2.1-fold over IgG, and Pol~II enrichment decreased from 11.3-fold to 3.4-fold, representing approximately 4-fold and 3-fold reductions, respectively.

\begin{table}[ht]
\centering
\caption{ChIP-qPCR at lncRNA-producing loci (8~h post-mixing).}
\label{tab:chip}
\begin{tabular}{lcc}
\toprule
\textbf{Factor} & \textbf{WT (fold over IgG)} & \textbf{Ema2-KO (fold over IgG)} \\
\midrule
Spt6 & 8.5 & 2.1 \\
RNA Pol II (Rpb1) & 11.3 & 3.4 \\
\bottomrule
\end{tabular}
\end{table}

% ============================================================
\section{Discussion}

Our findings establish Ema2 as a SUMO E3 ligase that plays a central role in the lncRNA transcription program required for scnRNA-directed DNA elimination in \textit{Tetrahymena}. We propose the following mechanistic model: Ema2 recruits the E2 conjugating enzyme Ubc9 through its RING domain to SUMOylate the transcription elongation factor Spt6. SUMOylated Spt6, in turn, promotes the stable association of both itself and RNA Polymerase~II with chromatin at somatic lncRNA-producing loci. This facilitates the robust transcription of lncRNAs that serve as templates for scnRNA comparison, enabling TDSD and ultimately DNA elimination.

This model is supported by several lines of evidence. First, Ema2 knockout nearly abolishes both lncRNA accumulation (6--10-fold reduction) and progeny viability (from 95\% to 5\%). Second, Ema2 interacts with Ubc9 in a RING domain-dependent manner, consistent with canonical SUMO E3 ligase function. Third, Spt6 SUMOylation is reduced to 15\% of wild-type levels in Ema2-KO cells. Fourth, both Spt6 and Pol~II chromatin occupancy at lncRNA loci are reduced 3--4-fold in the absence of Ema2.

The identification of SUMOylation as a regulatory mechanism for lncRNA transcription during genome rearrangement adds a new dimension to our understanding of programmed DNA elimination. While previous work has focused on the downstream effectors of scnRNA-guided heterochromatin formation, our results highlight the upstream transcriptional step as a regulated and essential component of the pathway.

An important question for future work is whether Spt6 is the sole relevant substrate of Ema2, or whether additional targets contribute to lncRNA transcription regulation. It will also be informative to determine whether analogous SUMOylation-dependent mechanisms operate in other ciliates that undergo programmed genome rearrangement, such as \textit{Oxytricha} and \textit{Paramecium}.

% ============================================================
\section{Conclusion}

We have identified Ema2 as a SUMO E3 ligase essential for long non-coding RNA transcription during programmed DNA elimination in \textit{Tetrahymena thermophila}. Ema2-mediated SUMOylation of Spt6 promotes the chromatin association of the transcription machinery at somatic lncRNA-producing loci, providing the lncRNA substrates required for scnRNA comparison and target-directed small RNA degradation. This work reveals SUMOylation as a key regulatory mechanism in genome rearrangement and connects post-translational modification signaling to the small RNA-directed epigenetic pathway.

% ============================================================
\bibliographystyle{plainnat}
\bibliography{references}

\end{document}
